\documentclass[10pt]{mypackage}

\usepackage[uprightgreeks,noDcommand]{kpfonts}
\renewcommand{\coloneq}{\coloneqq}
\renewcommand{\eqcolon}{\eqqcolon}
%\usepackage{dsfont}
%\renewcommand{\mathbb}{\mathds}
%\usepackage{homework}
\usepackage{notes}

\usepackage[ backend=bibtex, style = alphabetic, sorting=ynt ]{biblatex}
\addbibresource{../all_references.bib}

\usepackage{parskip}
\frenchspacing

\fancyhf{}
\fancyhead[R]{Avinash Iyer}
\fancyhead[L]{Random Walks and Varopoulos's Theorem}
\fancyfoot[C]{\thepage}

\setcounter{secnumdepth}{0}

\begin{document}
\RaggedRight
\section{Markov Chains}%
\begin{definition}
  The \textit{simple random walk} on $\Z^d$ is given by a walker starting at $\left( 0,\dots,0 \right)\in \Z^d$, then at time $t$, moves from position 
  \begin{align*}
    x(t) = \left( x_1,\dots,x_i,x_{i+1},\dots,x_d \right)
  \end{align*}
  to position
  \begin{align*}
    x\left( t+1 \right) &= \left( x_1,\dots,x_i,x_{i+1},\dots,x_d \right)
  \end{align*}
  with probability $1/(2d)$ for all $1\leq i \leq d$ and every discrete time unit.
\end{definition}
Our question is somewhat simple: when, if ever, does the walker return to the origin? This is the substance of Polya's Theorem and its generalization, Varopoulos's Theorem.
\begin{definition}
  Let $X$ be a finite or countable set. Let $P = \left( p\left( x,y \right) \right)_{x,y\in X}$ be a \textit{stochastic matrix}, a matrix satisfying $p\left( x,y \right) \geq 0$ for all $x,y\in X$ and
  \begin{align*}
    \sum_{y\in X}p\left( x,y \right) = 1
  \end{align*}
  for all $x\in X$.

  We call $P$ the one-step transition probability matrix. The \textit{initial distribution} is a map $\pi\colon X\to \R$ such that $\pi(x)\geq 0$ for all $x\in X$ and $\sum_{x\in X}\pi(x) = 1$. The triple $\left( X,P,\pi \right)$ is called a \textit{Markov chain.}
\end{definition}
We may formalize this as follows. Let
\begin{align*}
  \Omega &= X^{\N}
\end{align*}
be the space of trajectories. For a fixed $y_0,\dots,y_m\in X$, we let
\begin{align*}
  C\left( y_0,\dots,y_m \right) &= \bigcap_{j=0}^{m}Z_j^{-1}\left( \set{y_j} \right)
\end{align*}
be the cylinder set with base tuple $\left( y_0,\dots,y_m \right)\in X^{m+1}$, where we let $Z_j\colon \Omega\to X$ be the projection map. Define a probability measure $ \mathbb{P}_{\pi} $ by
\begin{align*}
  \mathbb{P}_{\pi}\left( C\left( y_0,\dots,y_m \right) \right) &= \pi\left( y_0 \right)p\left( y_0,y_1 \right)\cdots p\left( y_{m-1}y_m \right).
\end{align*}
By the Kolmogorov Extension Theorem, we may uniquely extend $ \mathbb{P}_{\pi} $ to the $\sigma$-algebra $ \mathcal{B}\subseteq P\left( \Omega \right) $ generated by the cylinder sets. We will let $ \mathbb{P} $ be this probability measure.

The projection maps $Z_j$ are then $X$-valued random variables reflecting the position of the walker at time $j$. The projection maps have the special property known as memorylessness --- that is, 
\begin{align*}
  p\left( x,y \right) &= \mathbb{P}\left( Z_{n+1}=y|Z_n = x \right)\\
                      &= \mathbb{P}\left( Z_{n+1}=y|Z_n = x\cap Z_{n-1}=x_{n-1}\cap \dots\cap Z_0=x_0 \right).
\end{align*}
This is also known as the Markov property.

Conversely, it can be shown that if $\left( Z_n \right)_n$ is a sequence of $X$-valued random variables satisfying the Markov property, then there is a Markov chain $\left( X,P,\pi \right)$ that can be found by taking
\begin{align*}
  \pi(x) &= \mathbb{P}\left( Z_0 = x \right)
\end{align*}
and determining $p$ as follows: if there exists $n\in \N$ such that $P\left( Z_n=x \right) > 0$, set $n(x)$ to be the minimal such value, and else set $n(x) = \infty$; then, set
\begin{align*}
  p\left( x,y \right) &= \begin{cases}
    P\left( Z_{n(x)+1}=y | P\left( Z_{n(x) = x} \right) \right) & n(x) < \infty\\
    \Phi(y)& \text{else}
  \end{cases}
\end{align*}
for all $x,y\in X$, where $\Phi\colon X\to [0,1]$ is a map such that $\sum_{y\in X}\Phi(y) = 1$. If $n(x) < \infty$ for all $x\in X$, then this Markov chain is unique.

Note that $\Phi$ is, in a sense, an atomic probability measure. 
\nocite{topics_in_groups_and_geometry}
\printbibliography
\end{document}
